\chapter*{Antes de começar}
\addcontentsline{toc}{chapter}{Antes de começar}
\markboth{Antes de começar}{Antes de começar}

Este guia foi escrito do zero, em português, com explicações próprias, para você aprender econometria de corte transversal, dados em painel e séries temporais. Ele \emph{não} é uma tradução de nenhum dos livros. Os livros enviados são a base de leitura aprofundada. Este material é o caminho didático para chegar até eles.

\section*{Como o guia está organizado}
Cada capítulo segue a mesma ordem: a ideia em linguagem simples, a matemática essencial, um experimento numérico que confirma a teoria, as armadilhas mais comuns e exercícios com solução. As caixas coloridas têm função fixa:

\begin{intuicao}[{Intuição}]
Explica o \emph{porquê} antes da fórmula.
\end{intuicao}
\begin{atencao}[{Armadilha}]
Erros frequentes, inclusive os que aparecem em livros e em manuais de software.
\end{atencao}
\begin{simulacao}[{Conferido por simulação}]
Resultado que eu rodei no computador (Python) para ter certeza de que a afirmação está correta. Os números são de uma semente fixa, então variam um pouco se você repetir.
\end{simulacao}

\section*{Mapa dos livros: onde aprofundar cada tema}
\begin{center}
\small
\begin{tabular}{p{4.3cm}p{9.6cm}}
\toprule
\textbf{Tema} & \textbf{Onde ler mais} \\
\midrule
Probabilidade e álgebra linear de base & Stachurski (\emph{Econometric Theory}, notas de curso), partes I e II; Goldberger \\
Regressão, MQO, VI & Wooldridge (cap. 4 e 5 da 1ª edição); Hanck et al. (cap. 4 a 9 e 12) \\
Dados em painel & Wooldridge (cap. 10 e 11); Hanck et al. (cap. 10) \\
ARMA, ARIMA, Box--Jenkins & Cryer e Chan (cap. 4 a 9); Mills (cap. 3, 4 e 7); Hamilton (cap. 3 e 4) \\
Raízes unitárias e cointegração & Hamilton (cap. 17 a 19); Mills (cap. 14); Hanck et al. (cap. 14 e 16) \\
VAR & Lütkepohl (\emph{New Introduction to Multiple Time Series Analysis}); Hamilton (cap. 11); Mills (cap. 13) \\
Volatilidade (ARCH/GARCH) & Tsay (\emph{Analysis of Financial Time Series}); Mills (cap. 10); Hamilton (cap. 21) \\
Aprendizado de máquina & Hastie, Tibshirani e Friedman (\emph{Elements of Statistical Learning}); tutorial de Faloutsos et al. sobre previsão em larga escala; Faouzi e Janati (pacote \texttt{pyts}) \\
\bottomrule
\end{tabular}
\end{center}

\section*{Duas observações sobre os livros}
\textbf{Edição do Wooldridge.} O arquivo que você enviou é a 1ª edição de \emph{Econometric Analysis of Cross Section and Panel Data} (2002, com impressão de 2004). Existe uma 2ª edição, de 2010, bem mais extensa. Os conceitos que uso aqui (MQO, VI, efeitos fixos e aleatórios) estão nas duas, mas a numeração de capítulos e alguns tópicos avançados mudam. Para o trabalho aplicado atual, vale ter a 2ª edição.

\textbf{Sobre ``corrigir os erros dos livros''.} Eu não auditei milhares de páginas atrás de erratas, e não vou fingir que sim. O que fiz foi (i) conferir por simulação toda fórmula e todo número que aparece neste guia e (ii) destacar, na caixa \emph{Armadilha}, os pontos onde os livros divergem entre si ou onde um erro é comum. O mais importante deles é o sinal do coeficiente de médias móveis, que muda de livro para livro (capítulo~\ref{cap:est}). Também não li os livros por inteiro: usei os sumários para montar o mapa acima e conferi diretamente em Cryer--Chan e em Mills a convenção de sinal do MA. Se você encontrar algo que não bate com seu livro, confie na simulação e na derivação, não no meu texto.

\section*{Notação}
Vetores e matrizes em negrito quando for útil. $\E(\cdot)$ é esperança, $\Var$ variância, $\Cov$ covariância, $\plim$ limite em probabilidade. $\hat\beta$ é um estimador e $\beta$ o parâmetro verdadeiro. Em séries temporais, $B$ é o operador de defasagem: $B y_t = y_{t-1}$.
